\chapter{The operator surfaces}
\label{ch:operator}

Three screens that users never see. They are for whoever runs the system --- and
during development, that is you.

\section{Admin: the audit log}

\shot{13-admin.png}{The administration screen. Members, services, and the audit
log.}

The audit log is the part worth studying. Every action that changes anything is
recorded here: who did it, to what, and when.

\note{This is a \emph{second} append-only record, and it is not the same thing
as the incident timeline. The timeline answers ``what happened to this
incident?'' and is visible to everyone. The audit log answers ``what did this
\emph{person} do?'' across the whole organisation, and only an administrator can
read it. One is the story of an event; the other is the story of an actor.}

You saw it being written during sign-in, in Chapter~\ref{ch:login}:

\where{api/app/Http/Controllers/Api/V1/AuthController.php}
\begin{lstlisting}[language=PHP]
$this->audit->record('auth.logged_in', $user, $user);
\end{lstlisting}

Three arguments: what happened, who did it, and what it happened to. For a
sign-in the actor and the subject are the same person, which is why \fpath{\$user}
appears twice.

\section{Horizon: watching the queue}

\shot{14-horizon-dashboard.png}{Laravel Horizon at
\fpath{http://localhost:8080/horizon}. Jobs processed, jobs failed, throughput.}

Chapter~\ref{ch:arch} said the API writes a notification row and returns
immediately, leaving \fpath{horizon} to do the sending. This is where you watch
that happen.

It answers questions that are otherwise invisible:

\begin{itemize}[leftmargin=1.4cm]
\item Are jobs being processed at all, or is the queue silently stuck?
\item Is anything failing repeatedly?
\item How long is work waiting before it is picked up?
\end{itemize}

\warnbox{This dashboard was completely broken for most of this project's life,
and nothing detected it. The page returned HTTP 200. The server logs were clean.
The health checks passed. It simply rendered a blank dark rectangle, and the
only evidence anywhere was in the browser's developer console.
Chapter~\ref{ch:bugs-infra} is that investigation in full --- it is the single
best example in this book of a failure that no amount of reading the server logs
would ever have found.}

\section{Mailpit: catching email that must not escape}

\shot{15-mailpit.png}{Mailpit at \fpath{http://localhost:8025}. Every email the
system has sent, and not one of them delivered to a real person.}

In development, IncidentFlow sends its email to \fpath{mailpit} instead of a
real mail server. Mailpit accepts everything and delivers nothing, showing it in
a web inbox instead.

\note{This is not merely convenient. The seeded database contains addresses like
\fpath{commander@incidentflow.test}, but a realistic test dataset eventually
contains something that looks like a real address --- and a \fpath{sev1}
notification mails every responder in the organisation. Without a catcher,
running the test suite once could send mail to strangers. Mailpit makes that
mistake impossible rather than unlikely.}

\subsection{Watch the whole chain in one action}

This is the best single exercise in the book, because it exercises five
containers at once.

\begin{enumerate}[leftmargin=1.5cm]
\item Open \fpath{http://localhost:8025} in one tab. It should be empty or
  nearly so.
\item In another tab, sign in as \fpath{commander@incidentflow.test} and create
  an incident with severity \textbf{sev1}.
\item Watch the Mailpit tab.
\end{enumerate}

\expect{Email appears in Mailpit within a second or two, one per responder in
the organisation. You have just watched: the \textbf{api} write rows and commit;
\textbf{redis} carry the job; \textbf{horizon} pick it up; \textbf{mailpit}
receive the result --- and \textbf{postgres} holding the record throughout.}

\warnbox{If no mail arrives, the fault is almost always Horizon rather than the
mail path. Check \fpath{docker compose logs horizon}, and check the Horizon
dashboard for failed jobs. A notification row that stays \fpath{pending} means
the job was never queued or never ran; one marked \fpath{failed} means it ran
and threw.}

\subsection{Try the same thing at sev3}

Create a second incident, this time \textbf{sev3}. No email arrives.

That is the rule from Chapter~\ref{ch:what}: only \fpath{sev1} and \fpath{sev2}
notify the whole responder pool. You have now seen one enum value decide whether
a hundred people are interrupted --- which is the concrete reason changing
severity is a permission most roles do not have.

\section{Check your understanding}

\begin{itemize}[leftmargin=1.4cm]
\cbox What question does the audit log answer that the incident timeline does
  not?
\cbox Why does sending email through a queue make the API faster \emph{and}
  more reliable?
\cbox Why is a mail catcher a safety feature rather than a convenience?
\cbox A notification row is stuck at \fpath{pending}. What are the two
  possibilities?
\end{itemize}
